\chapter{Diseases of the Kidneys and Urinary System}
\label{ch:kidneys}
The kidneys are paired retroperitoneal organs responsible for filtering blood, excreting metabolic waste, maintaining fluid and electrolyte balance, regulating blood pressure (via the renin-angiotensin-aldosterone system), and producing erythropoietin and active vitamin D. Each kidney contains approximately one million nephrons. Diseases of the kidneys range from glomerular disorders and tubulointerstitial diseases to structural abnormalities and neoplasms.

%-------------------------------------------------------------
\section{Acute Kidney Injury}
%-------------------------------------------------------------

%-------------------------------------------------------------

%-------------------------------------------------------------

%-------------------------------------------------------------

%-------------------------------------------------------------

%-------------------------------------------------------------

\label{sec:section:Acute Kidney Injury}
%-------------------------------------------------------------%-------------------------------------------------------------

\subsection{Acute Kidney Injury}
\label{sec:Acute Kidney Injury}
Acute kidney injury (AKI) is a rapid decline in kidney function over hours to days, resulting in accumulation of nitrogenous waste products (urea, creatinine) and inability to maintain fluid, electrolyte, and acid-base homeostasis. The KDIGO criteria define AKI as an increase in serum creatinine $\geq$0.3 mg/dL within 48 hours, increase in creatinine to $\geq$1.5 times baseline within 7 days, or urine volume $<$0.5 mL/kg/h for 6 hours. Causes are classified as prerenal (decreased renal perfusion: dehydration, heart failure, sepsis, NSAIDs), intrinsic (acute tubular necrosis from ischemia or nephrotoxins, glomerulonephritis, interstitial nephritis, thrombotic microangiopathy), and postrenal (obstruction: kidney stones, enlarged prostate, pelvic malignancy). The condition results from the specific underlying cause leading to reduced GFR. Clinical features include oliguria or anuria, fluid overload, electrolyte disturbances, and metabolic acidosis. Diagnosis is based on serum creatinine elevation, urine output criteria, urinalysis, and imaging. Management involves addressing the underlying cause, fluid and electrolyte management, and renal replacement therapy (dialysis) for severe cases. Prognosis depends on severity and cause; prerenal azotemia is reversible with restoration of perfusion; ATN typically recovers over days to weeks. Mortality is significant, particularly in hospitalized patients.

%-------------------------------------------------------------
\section{Chronic Kidney Disease}
%-------------------------------------------------------------

%-------------------------------------------------------------

%-------------------------------------------------------------

%-------------------------------------------------------------

%-------------------------------------------------------------

%-------------------------------------------------------------

\label{sec:section:Chronic Kidney Disease}
%-------------------------------------------------------------%-------------------------------------------------------------

\subsection{Chronic Kidney Disease}
\label{sec:Chronic Kidney Disease}
Chronic kidney disease (CKD) is abnormalities of kidney structure or function present for more than 3 months, with health implications. The most common causes are diabetes mellitus (leading cause in developed countries) and hypertension; other causes include glomerulonephritis, polycystic kidney disease, and chronic obstruction. The condition results from progressive loss of nephrons leading to declining GFR. CKD is staged by the glomerular filtration rate (GFR): Stage 1 (GFR $\geq$90 with kidney damage markers), Stage 2 (GFR 60--89), Stage 3a (GFR 45--59), Stage 3b (GFR 30--44), Stage 4 (GFR 15--29), and Stage 5 (GFR $<$15, end-stage renal disease, ESRD). Clinical features are often absent in early stages; advanced disease includes anemia (decreased erythropoietin), bone mineral disease (secondary hyperparathyroidism, vitamin D deficiency, hyperphosphatemia), metabolic acidosis, hyperkalemia, volume overload, and cardiovascular disease. Diagnosis is based on GFR estimation and markers of kidney damage (albuminuria, abnormal imaging, abnormal histology). Management includes BP control (target $<$130/80 with ACE inhibitors/ARBs), glycemic control, SGLT2 inhibitors (dapagliflozin, empagliflozin), dietary modification, and avoidance of nephrotoxins; renal replacement therapy is required for Stage 5 CKD. Prognosis depends on stage and progression rate; cardiovascular disease is the leading cause of death in CKD.

\subsection{End-Stage Renal Disease}
\label{sec:End-Stage Renal Disease}
End-stage renal disease (ESRD) is the final stage of CKD (GFR $<$15 mL/min/1.73 m$^2$) where the kidneys can no longer sustain life without renal replacement therapy. The condition results from progressive nephron loss from any cause of CKD. Clinical features include uremic syndrome: nausea, vomiting, anorexia, pruritus, pericarditis, encephalopathy, peripheral neuropathy, and restless legs syndrome. Diagnosis is based on GFR criteria and clinical manifestations. Management involves dialysis (hemodialysis or peritoneal dialysis) or kidney transplantation; kidney transplantation is the treatment of choice for eligible patients, offering superior survival and quality of life. Post-transplant immunosuppression (tacrolimus, mycophenolate, prednisone) carries risks of infection and malignancy. Prognosis is variable; living donor transplantation provides the best outcomes, while dialysis has a 5-year survival of approximately 40--50\%.

\subsection{Renal Osteodystrophy}
\label{sec:Renal Osteodystrophy}
Renal osteodystrophy is a spectrum of bone disorders that develop in patients with chronic kidney disease (CKD), resulting from disturbances in mineral metabolism, affecting virtually all patients with CKD stage 5. The condition results from phosphate retention, decreased renal production of active vitamin D (calcitriol), hypocalcemia, and secondary hyperparathyroidism. Bone histology subtypes include osteitis fibrosa cystica (high-turnover, from severe hyperparathyroidism), adynamic bone disease (low-turnover, from oversuppression of PTH with calcitriol and calcium-based binders, common in diabetics and the elderly), osteomalacia (defective mineralization, from vitamin D deficiency or aluminum toxicity), and mixed uremic osteodystrophy. Clinical features include bone pain (most commonly in lower back, hips, and legs), proximal muscle weakness, fractures, skeletal deformities, growth retardation in children, and extraosseous calcification (vascular calcification, calciphylaxis). Diagnosis is based on elevated PTH, hyperphosphatemia, hypocalcemia or hypercalcemia, low vitamin D levels, and bone biopsy (gold standard but rarely performed). Management involves phosphate restriction (dietary), phosphate binders (calcium-based, sevelamer, lanthanum carbonate, iron-based binders), active vitamin D analogs (calcitriol, paricalcitol, doxercalciferol), calcimimetics (cinacalcet, etelcalcetide), and optimization of dialysis. Parathyroidectomy for refractory hyperparathyroidism. Prognosis is improved by early recognition and treatment; renal osteodystrophy improves after successful kidney transplantation.

%-------------------------------------------------------------
\section{Glomerular Diseases}
%-------------------------------------------------------------

%-------------------------------------------------------------

%-------------------------------------------------------------

%-------------------------------------------------------------

%-------------------------------------------------------------

%-------------------------------------------------------------

\label{sec:Glomerular Diseases}
%-------------------------------------------------------------%-------------------------------------------------------------

\subsection{Diabetic Nephropathy}
\label{sec:Diabetic Nephropathy}
Diabetic nephropathy is the most common cause of ESRD worldwide, affecting 30--40\% of patients with diabetes mellitus (both T1DM and T2DM). The condition results from hyperglycemia-induced metabolic and hemodynamic changes leading to glomerular hyperfiltration, mesangial expansion, basement membrane thickening, and ultimately nodular glomerulosclerosis (Kimmelstiel-Wilson nodules). Clinical progression includes initial hyperfiltration (increased GFR), microalbuminuria (moderately increased albuminuria), macroalbuminuria (heavily increased albuminuria), declining GFR, and ESRD. Diagnosis is based on persistent albuminuria in the setting of diabetes. Management includes optimal glycemic control, RAAS blockade (ACE inhibitors or ARBs), SGLT2 inhibitors, BP control (target $<$130/80), and finerenone (non-steroidal mineralocorticoid receptor antagonist). Prognosis is improved with early intervention and control of risk factors.

\subsection{Focal Segmental Glomerulosclerosis}
\label{sec:Focal Segmental Glomerulosclerosis}
Focal segmental glomerulosclerosis (FSGS) is a pattern of glomerular injury characterized by segmental sclerosis in some (focal) glomeruli. It is a common cause of nephrotic syndrome in adults. The condition may be primary (idiopathic) or secondary (to obesity, HIV infection, sickle cell disease, drug abuse, or adaptive hyperfiltration); podocyte injury is the central pathological event. Clinical features include proteinuria (often nephrotic-range), hematuria, hypertension, and progressive renal insufficiency. Diagnosis is by renal biopsy, which must include sufficient glomeruli to demonstrate the focal pattern. Management of primary FSGS includes corticosteroids (6 months), calcineurin inhibitors for steroid-resistant cases, and rituximab; secondary FSGS requires treatment of the underlying cause. Prognosis is variable; APOL1 gene variants are associated with increased risk, particularly in individuals of African ancestry.

\subsection{IgA Nephropathy}
\label{sec:IgA Nephropathy}
IgA nephropathy (Berger's disease) is the most common primary glomerulonephritis worldwide, with higher prevalence in East Asia and Southern Europe. The condition results from deposition of galactose-deficient IgA1 in the mesangium, triggering an immune response, mesangial proliferation, and glomerular injury. Clinical features include episodic gross hematuria (often concurrent with upper respiratory infections---"synpharyngitic hematuria"), microscopic hematuria, or proteinuria; some present with nephrotic syndrome or rapidly progressive glomerulonephritis. Diagnosis requires renal biopsy showing mesangial IgA deposition on immunofluorescence; the Oxford Classification (MEST-C scoring) guides prognosis. Management depends on severity: RAAS blockade for proteinuria, SGLT2 inhibitors, and immunosuppression (corticosteroids) for high-risk patients; supportive care includes BP control and smoking cessation. Neflgentan (sodium-glucose cotransporter 2 inhibitor) is an emerging therapy. Prognosis is variable; up to 30\% progress to ESRD within 20 years.

\subsection{Membranous Nephropathy}
\label{sec:Membranous Nephropathy}
Membranous nephropathy (MN) is the most common cause of nephrotic syndrome in non-diabetic adults. The condition results from subepithelial immune complex deposition and thickening of the glomerular basement membrane; approximately 70--80\% of primary MN is caused by autoantibodies against the phospholipase A2 receptor (PLA2R) on podocytes. Secondary causes include malignancy (lung, GI), infections (hepatitis B), autoimmune diseases (lupus), and drugs. Clinical features include nephrotic-range proteinuria, edema, and hypoalbuminemia. Diagnosis is based on renal biopsy with PLA2R antibody serology. Management of primary MN includes conservative management (RAAS blockade) for mild disease and immunosuppression (rituximab, cyclophosphamide, calcineurin inhibitors) for moderate-to-severe or progressive disease. Prognosis is variable; anti-PLA2R antibody levels monitor disease activity and response to treatment.

\subsection{Minimal Change Disease}
\label{sec:Minimal Change Disease}
Minimal change disease (MCD) is the most common cause of nephrotic syndrome in children (70--90\% of cases) and a significant cause in adults (10--15\%). The condition results from podocyte injury with foot process effacement on electron microscopy, with normal light microscopy and immunofluorescence (hence "minimal change"). Most cases are idiopathic, though MCD can be associated with NSAIDs, lymphoma, and infections. Clinical features include sudden onset of nephrotic-range proteinuria (>3.5 g/day), edema, hypoalbuminemia, hyperlipidemia, and lipiduria. Diagnosis is based on renal biopsy demonstrating podocyte foot process effacement on electron microscopy. Management is with corticosteroids (prednisone), to which 90\% of patients respond (steroid-sensitive); frequent relapses or steroid dependence may require second-line agents (calcineurin inhibitors, cyclophosphamide, rituximab). Prognosis is excellent; most patients achieve remission with corticosteroids.

%-------------------------------------------------------------
\section{Kidney Stones}
%-------------------------------------------------------------

%-------------------------------------------------------------

Kidney stones, or nephrolithiasis, result from super-saturation, nucleation, and aggregation of crystalline mineral compounds within the renal pelvis and calyces. Obstruction of the ureteral lumen by displaced calculi causes severe flank pain (renal colic), hematuria, and risk of obstructive uropathy or pyelonephritis. Management includes pain control, medical expulsive therapy, extracorporeal shock wave lithotripsy, ureteroscopy, and preventive dietary optimization.

%-------------------------------------------------------------

%-------------------------------------------------------------

%-------------------------------------------------------------

%-------------------------------------------------------------

\label{sec:Kidney Stones}
%-------------------------------------------------------------%-------------------------------------------------------------

\subsection{Nephrolithiasis}
\label{sec:Nephrolithiasis}
Nephrolithiasis (kidney stones) affects approximately 10\% of the population, with higher prevalence in males and hot climates. Stone types include calcium oxalate (most common, 70--80\%), calcium phosphate, uric acid (associated with hyperuricemia, low urine pH, and metabolic syndrome), struvite (magnesium ammonium phosphate, associated with urease-producing organisms in chronic UTI---"staghorn calculi"), and cystine (rare, hereditary cystinuria). The condition results from supersaturation of stone-forming salts in urine. Risk factors include low fluid intake, high dietary sodium and animal protein, hypercalciuria, hyperoxaluria, hypocitraturia, and hyperuricosuria. Clinical features include renal colic (severe, colicky flank pain radiating to the groin), hematuria, nausea, and urinary urgency. Diagnosis is based on CT abdomen/pelvis without contrast (gold standard for acute evaluation). Management of acute renal colic includes NSAIDs, opioids, antiemetics, and alpha-blockers (medical expulsive therapy). Prevention includes increased fluid intake (>2.5 L/day), dietary modifications, and pharmacotherapy (thiazide diuretics for hypercalciuria, allopurinol for uric acid stones, potassium citrate for hypocitraturia). Prognosis is good with preventive measures; recurrence rates are high without prophylaxis.

%-------------------------------------------------------------
\section{Lower Urinary Tract Disorders}
%-------------------------------------------------------------

%-------------------------------------------------------------

%-------------------------------------------------------------

%-------------------------------------------------------------

%-------------------------------------------------------------

%-------------------------------------------------------------

\label{sec:Lower Urinary Tract Disorders}
%-------------------------------------------------------------%-------------------------------------------------------------

\subsection{Urinary Incontinence}
\label{sec:Urinary Incontinence}
Urinary incontinence is the involuntary loss of urine, affecting approximately 25--45\% of women and 10--25\% of men, with prevalence increasing with age. Types include stress urinary incontinence (SUI, leakage with coughing, sneezing, laughing, or exercise, caused by urethral hypermobility or intrinsic sphincter deficiency), urge urinary incontinence (UUI, involuntary loss of urine preceded by a sudden, strong urge, associated with detrusor overactivity), mixed urinary incontinence (combination of SUI and UUI), and overflow incontinence (continuous dribbling from an overdistended bladder due to bladder outlet obstruction or detrusor underactivity). The condition results from dysfunction of the pelvic floor, urethral sphincter, or detrusor muscle. Risk factors for SUI include vaginal delivery, menopause, obesity, and chronic cough; UUI may be idiopathic or secondary to neurogenic bladder, bladder stones, infection, or malignancy. Clinical features depend on the type of incontinence. Diagnosis is based on history, bladder diary, physical examination (cough stress test, pelvic organ prolapse assessment), urinalysis, and post-void residual volume measurement; urodynamic studies are indicated for refractory cases. Management of SUI includes pelvic floor muscle training (Kegel exercises, first-line), lifestyle modification, vaginal pessaries, and surgical options (mid-urethral sling procedures, colposuspension); management of UUI includes bladder retraining, timed voiding, antimuscarinics, beta-3 agonists (mirabegron, vibegron), and onabotulinumtoxinA for refractory cases; overflow incontinence is managed by treating the underlying obstruction or optimizing bladder emptying. Prognosis is generally favorable with appropriate management tailored to the type of incontinence.

%-------------------------------------------------------------
\section{Polycystic Kidney Disease}
%-------------------------------------------------------------

%-------------------------------------------------------------

%-------------------------------------------------------------

%-------------------------------------------------------------

%-------------------------------------------------------------

%-------------------------------------------------------------

\label{sec:Polycystic Kidney Disease}
%-------------------------------------------------------------%-------------------------------------------------------------

\subsection{Autosomal Dominant Polycystic Kidney Disease}
\label{sec:Autosomal Dominant Polycystic Kidney Disease}
Autosomal dominant polycystic kidney disease (ADPKD) is the most common hereditary kidney disease, affecting 1 in 400--1000 individuals. It is caused by mutations in PKD1 (chromosome 16, 85\% of cases) or PKD2 (chromosome 4, 15\% of cases). The condition results from defective polycystin proteins leading to progressive bilateral renal cyst formation, massive kidney enlargement, declining GFR, and ESRD typically by the fifth to seventh decade. Clinical features include flank pain, hematuria, hypertension, and abdominal distension; extrarenal manifestations include hepatic cysts (most common), intracranial berry aneurysms (5--10\%), cardiac valvular abnormalities, and colonic diverticulosis. Diagnosis is based on imaging (CT or MRI showing $\geq$3 cysts per kidney) and genetic testing. Management includes BP control (ACE inhibitors/ARBs), salt restriction, adequate hydration, and tolvaptan (vasopressin V2 receptor antagonist) to slow cyst growth and kidney function decline. Prognosis is variable; screening of first-degree relatives is recommended.

%-------------------------------------------------------------
\section{Renal and Urologic Neoplasms}
%-------------------------------------------------------------

%-------------------------------------------------------------

%-------------------------------------------------------------

%-------------------------------------------------------------

%-------------------------------------------------------------

%-------------------------------------------------------------

\label{sec:Renal and Urologic Neoplasms}
%-------------------------------------------------------------%-------------------------------------------------------------

\subsection{Bladder Cancer}
\label{sec:Bladder Cancer}
Bladder cancer is the fourth most common cancer in men, with 550,000 new cases/year worldwide. Risk factors include smoking (50\% of cases), occupational exposure (aromatic amines, aniline dyes), chronic cystitis (squamous cell carcinoma), \textit{Schistosoma haematobium} infection (squamous cell in endemic areas), cyclophosphamide, pelvic radiation, and family history. Urothelial (transitional cell) carcinoma accounts for >90\%; squamous cell carcinoma (3--5\%) and adenocarcinoma (1--2\%) are less common. The condition results from multifocal field cancerization with mutations in FGFR3 (low-grade) and p53/Rb (high-grade, CIS). Clinical features include gross, painless hematuria (most common), microscopic hematuria, and irritative voiding symptoms with CIS. Diagnosis is based on cystoscopy (gold standard) with biopsy/resection, urine cytology, FISH (UroVysion), and CT urography. Staging distinguishes non-muscle-invasive (Ta, T1, CIS) from muscle-invasive (T2+) disease. Management of NMIBC includes TURBT with risk-adapted intravesical therapy (mitomycin C, BCG); BCG-unresponsive NMIBC requires radical cystectomy or pembrolizumab. MIBC is treated with neoadjuvant cisplatin-based chemotherapy followed by radical cystectomy with urinary diversion; trimodal therapy is an option for patients unfit for cystectomy. Metastatic disease is treated with platinum-based chemotherapy, checkpoint inhibitors, FGFR inhibitors (erdafitinib for FGFR3-mutant), and antibody-drug conjugates (enfortumab vedotin, sacituzumab govitecan). Prognosis varies by stage: 5-year survival is 96\% for Ta, 70--90\% for T1, 50--60\% for T2, 30--50\% for T3, and <15\% for metastatic.

\subsection{Kidney Cancer (Renal Cell Carcinoma)}
\label{sec:Kidney Cancer}
Renal cell carcinoma (RCC) accounts for 2--3\% of adult malignancies (400,000 cases/year worldwide). Risk factors include smoking, obesity, hypertension, acquired cystic kidney disease, and hereditary syndromes (VHL, hereditary papillary RCC, Birt-Hogg-Dubé, tuberous sclerosis, hereditary leiomyomatosis). Clear cell RCC (75\%) is the most common subtype, followed by papillary (15\%), chromophobe (5\%), and collecting duct carcinoma. The condition results from VHL inactivation leading to HIF accumulation, VEGF overexpression, and angiogenesis in clear cell RCC. Clinical features include the triad of hematuria, flank pain, and palpable abdominal mass (now uncommon; 60\% are found incidentally on imaging); paraneoplastic syndromes include erythrocytosis (EPO), hypercalcemia (PTHrP), hypertension (renin), and Stauffer syndrome (reversible liver dysfunction). Diagnosis is based on CT abdomen with contrast (heterogeneous, enhancing renal mass), MRI for cystic/complex lesions, and renal mass biopsy for small tumors; staging uses the TNM system. Management includes partial nephrectomy (nephron-sparing for T1a, T1b tumors), radical nephrectomy for larger tumors, and ablative therapies for small tumors in elderly patients; advanced disease is treated with tyrosine kinase inhibitors (sunitinib, pazopanib, cabozantinib), immunotherapy (nivolumab + ipilimumab, pembrolizumab + axitinib), mTOR inhibitors (everolimus, temsirolimus), and HIF-2$\alpha$ inhibitors (belzutifan for VHL-associated RCC). Prognosis varies by stage: 5-year survival is 93\% for localized, 72\% regional, and 14\% metastatic.

%-------------------------------------------------------------
\section{Renal Tubular Disorders}
%-------------------------------------------------------------

%-------------------------------------------------------------

%-------------------------------------------------------------

%-------------------------------------------------------------

%-------------------------------------------------------------

%-------------------------------------------------------------

\label{sec:Renal Tubular Disorders}
%-------------------------------------------------------------%-------------------------------------------------------------

\subsection{Fanconi Syndrome}
\label{sec:Fanconi Syndrome}
Fanconi syndrome is generalized dysfunction of the proximal renal tubule, leading to impaired reabsorption of glucose, amino acids, phosphate, bicarbonate, uric acid, and other solutes. Causes include hereditary disorders (cystinosis, Wilson's disease, galactosemia), drugs (ifosfamide, outdated tetracycline, aminoglycosides), multiple myeloma, and heavy metal poisoning. The condition results from impaired proximal tubular transport mechanisms. Clinical features include metabolic acidosis (proximal RTA), hypophosphatemic rickets/osteomalacia, polyuria, and growth retardation in children. Diagnosis is based on laboratory findings of generalized proximal tubular dysfunction. Management addresses the underlying cause and provides supplementation of lost solutes (phosphate, vitamin D, bicarbonate). Prognosis depends on the underlying cause.

\subsection{Renal Tubular Acidosis}
\label{sec:Renal Tubular Acidosis}
Renal tubular acidosis (RTA) is a group of disorders in which the kidneys fail to appropriately acidify the urine despite normal or near-normal GFR. Type 1 (distal) RTA results from impaired hydrogen ion secretion in the collecting duct, causing hyperchloremic metabolic acidosis with inappropriately alkaline urine (pH >5.5); complications include nephrocalcinosis and rickets. Type 2 (proximal) RTA involves impaired bicarbonate reabsorption in the proximal tubule, causing metabolic acidosis with appropriately acidic urine and bicarbonaturia; it is often part of Fanconi syndrome. Type 4 RTA is caused by aldosterone deficiency or resistance, leading to hyperkalemic metabolic acidosis. The condition results from the specific tubular transport defect. Clinical features vary by type and include metabolic acidosis, growth retardation, and nephrocalcinosis. Diagnosis is based on serum electrolytes, urine pH, and acid-base testing. Management includes oral bicarbonate supplementation (Types 1 and 2) and management of hyperkalemia and aldosterone deficiency (Type 4). Prognosis is favorable with appropriate supplementation.

%-------------------------------------------------------------
\section{Tubulointerstitial Diseases}
%-------------------------------------------------------------

%-------------------------------------------------------------

%-------------------------------------------------------------

%-------------------------------------------------------------

%-------------------------------------------------------------

%-------------------------------------------------------------

\label{sec:Tubulointerstitial Diseases}
%-------------------------------------------------------------%-------------------------------------------------------------

\subsection{Acute Interstitial Nephritis}
\label{sec:Acute Interstitial Nephritis}
Acute interstitial nephritis (AIN) is an immune-mediated inflammatory condition of the renal tubules and interstitium, a common cause of AKI in hospitalized patients. The most common cause is drug reaction (antibiotics [penicillins, cephalosporins, fluoroquinolones], NSAIDs, PPIs, diuretics); other causes include infections, autoimmune diseases, and idiopathic. The condition results from a hypersensitivity reaction to the offending agent. Clinical features may include fever, rash, eosinophilia, and eosinophiluria (the classic triad of fever, rash, and arthralgias is uncommon); AKI with sterile pyuria and mild proteinuria should raise suspicion. Diagnosis requires renal biopsy showing lymphocytic infiltration, edema, and tubulitis. Management involves discontinuation of the offending drug; corticosteroids are used for severe or delayed recovery. Prognosis is generally favorable with early identification and drug withdrawal.

\subsection{Chronic Tubulointerstitial Nephritis}
\label{sec:Chronic Tubulointerstitial Nephritis}
Chronic tubulointerstitial nephritis is gradual tubular atrophy and interstitial fibrosis. Causes include chronic analgesic use (NSAIDs, combined analgesics), lithium therapy, heavy metals (lead, cadmium), hyperuricemia, sarcoidosis, and chronic infections. The condition results from prolonged exposure to nephrotoxic agents or chronic inflammation. Clinical features include slowly progressive CKD, often with inability to concentrate urine (polyuria, nocturia) and mild proteinuria. Diagnosis is by renal biopsy showing tubular atrophy, interstitial fibrosis, and chronic inflammatory infiltrate. Management involves removal of the offending agent and management of CKD complications. Prognosis is variable; progression to ESRD may occur with continued exposure.

%-------------------------------------------------------------
\section{Urinary Tract Infections}
%-------------------------------------------------------------

%-------------------------------------------------------------

%-------------------------------------------------------------

%-------------------------------------------------------------

%-------------------------------------------------------------

%-------------------------------------------------------------

\label{sec:Urinary Tract Infections}
%-------------------------------------------------------------%-------------------------------------------------------------

\subsection{Cystitis}
\label{sec:Cystitis}
Cystitis (lower urinary tract infection) is one of the most common bacterial infections, particularly in women (50--60\% of women experience at least one episode). \textit{Escherichia coli} is the causative organism in 80--90\% of cases. The condition results from bacterial colonization and inflammation of the bladder mucosa. Clinical features include dysuria, urinary frequency, urgency, suprapubic pain, and hematuria. Diagnosis is based on urinalysis and urine culture. Management of uncomplicated cystitis in non-pregnant women includes nitrofurantoin, trimethoprim-sulfamethoxazole, or fosfomycin; recurrent cystitis may require prophylactic antibiotics, behavioral modifications, and cranberry products; complicated cystitis (in males, pregnant women, or with structural abnormalities) requires broader-spectrum antibiotics and evaluation for underlying causes. Prognosis is excellent with appropriate antibiotic therapy.

\subsection{Nephritic Syndrome}
\label{sec:Nephritic Syndrome}
Nephritic syndrome is a clinical complex resulting from acute glomerular inflammation characterized by leukocyte infiltration, cellular proliferation, and reactive glomerular capillary wall injury. The condition results from immune complex deposition or antibody binding within the glomerulus, causing a decrease in glomerular filtration rate (GFR) and breakdown of the capillary barrier to red blood cells and protein. Etiologies include poststreptococcal glomerulonephritis (PSGN), IgA nephropathy (Berger disease), lupus nephritis, and anti-GBM disease (Goodpasture syndrome). Clinical features include the classic triad of hematuria ('cola-colored' or dysmorphic RBCs and RBC casts), hypertension (due to fluid retention), and oliguria with mild-to-moderate proteinuria ($<3.5$ g/day) and facial/periorbital edema. Diagnosis is supported by urinalysis showing hematuria and proteinuria, elevated serum creatinine, low serum complement levels (in PSGN and lupus), and definitive renal biopsy. Management involves fluid and sodium restriction, antihypertensive therapy (diuretics, ACE inhibitors), and targeted immunosuppressants (corticosteroids, cyclophosphamide) for rapidly progressive cases. Prognosis depends on underlying histology.

\subsection{Nephrotic Syndrome}
\label{sec:Nephrotic Syndrome}
Nephrotic syndrome is a clinical complex resulting from severe glomerular podocyte effacement and disruption of the charge and size barrier of the glomerular basement membrane. The condition results from primary renal podocytopathies or systemic disease damage to the filtration barrier. Primary etiologies include minimal change disease (most common in children), focal segmental glomerulosclerosis (FSGS, most common in African Americans), and membranous nephropathy (most common in Caucasian adults); secondary etiologies include diabetic nephropathy, amyloidosis, and lupus nephritis. Clinical features include the hallmark tetrad: heavy proteinuria ($>3.5$ g/24 hours or spot urine protein-to-creatinine ratio $>3.5$), hypoalbuminemia ($<3.0$ g/dL), generalized peripheral and periorbital edema (anasarca due to decreased oncotic pressure and secondary sodium retention), and hyperlipidemia with lipiduria ('oval fat bodies' and Maltese cross appearance under polarized light). Complications include deep vein thrombosis/pulmonary embolism (due to urinary loss of antithrombin III) and infection (loss of immunoglobulins). Diagnosis is by 24-hour urine collection or spot ratio, serum chemistries, and renal biopsy. Management includes ACE inhibitors/ARBs (to reduce intraglomerular pressure), loop diuretics, statins, anticoagulation when indicated, and systemic corticosteroids or immunosuppressive agents. Prognosis varies by underlying etiology.

\subsection{Pyelonephritis}
\label{sec:Pyelonephritis}
Pyelonephritis is an infection of the renal parenchyma and renal pelvis, representing an upper urinary tract infection. Acute pyelonephritis typically results from ascending infection from the lower urinary tract, most commonly caused by \textit{E.~coli}. Risk factors include vesicoureteral reflux, urinary obstruction, diabetes, and immunosuppression. Clinical features include high fever, chills, flank pain, nausea, and vomiting; CVA tenderness is present on examination. Diagnosis is based on urinalysis (pyuria, bacteriuria, hematuria) and urine culture; blood cultures are indicated for severe or complicated cases. Management of mild cases is outpatient oral fluoroquinolones or initial IV ceftriaxone followed by oral step-down therapy; severe cases require hospitalization and IV antibiotics. Prognosis is favorable with appropriate antibiotic therapy; chronic pyelonephritis results from recurrent infections and may lead to renal scarring and CKD.

\subsection{Urosepsis}
\label{sec:Urosepsis}
Urosepsis is sepsis originating from the urinary tract, accounting for approximately 25\% of all sepsis cases. It is more common in the elderly, immunocompromised, and those with urological abnormalities. Common pathogens include \textit{E.~coli}, \textit{Klebsiella}, \textit{Pseudomonas}, and \textit{Enterococcus}. The condition results from a dysregulated host response to urinary tract infection leading to life-threatening organ dysfunction. Clinical features include fever, hypotension, altered mental status, and signs of organ dysfunction. Diagnosis is based on clinical criteria for sepsis with urinary source identified. Management follows the Surviving Sepsis Campaign guidelines: early recognition, broad-spectrum IV antibiotics within 1 hour, aggressive fluid resuscitation, vasopressors for septic shock, source control (e.g., decompression of obstructed urinary tract), and organ support. Prognosis is guarded; mortality ranges from 20--40\% depending on severity.

